%==============================================================================
% Poster bachelorproef
%==============================================================================
% Gebaseerd op document class `a0poster' door Gerlinde Kettl en Matthias Weiser
% Aangepast voor gebruik aan HOGENT door Jens Buysse en Bert Van Vreckem

\documentclass[a0,portrait]{hogent-poster}

% Info over de opleiding
\course{Bachelorproef}
\studyprogramme{toegepaste informatica}
\academicyear{2025-2026}
\institution{Hogeschool Gent, Valentin Vaerwyckweg 1, 9000 Gent}

% Info over de bachelorproef
\title{Automatische netwerkdetectie in NetBox}
\subtitle{Een hybride discoverypipeline voor een betrouwbare single source of truth bij Citymesh}
\author{Laurens Lecocq}
\email{laurens.lecocq@student.hogent.be}
\supervisor{Martijn Saelens}
\cosupervisor{Tom Pirlet (Citymesh)}

\specialisation{Systeem- en Netwerkbeheer}
\keywords{NetBox Discovery, Diode, Orb, LLDP, SNMP, NetworkTool, single source of truth}
% \projectrepo{https://github.com/user/repo}

\begin{document}

\maketitle

\begin{abstract}
In dynamische netwerkomgevingen veroudert documentatie snel, waardoor een betrouwbare \emph{single source of truth} (SSOT) moeilijk te behouden is. Deze bachelorproef onderzoekt hoe automatische netwerkdetectie de actualiteit en juistheid van netwerkdocumentatie in NetBox kan verbeteren op WiFi-/AP-sites bij Citymesh. De kern van de bijdrage is een \textbf{hybride discoverypipeline}: NetBox Discovery met Diode voedt de inventarislaag, terwijl \texttt{NetworkTool}, een eigen ontwikkeling, de LLDP-gebaseerde kabel- en topologielaag en de RUCKUS vSZ-koppeling afdekt. Een casestudy op vier sites toont dat deze aanpak de documentatiekwaliteit voor alle bereikbare apparatuur volledig herstelt.
\end{abstract}

\begin{multicols}{2}

%------------------------------------------------------------------------------
\section{Context en probleem}

Citymesh beheert sterk gedistribueerde, snel veranderende netwerkinfrastructuur. In zo'n omgeving raakt NetBox-documentatie snel verouderd wanneer wijzigingen niet consequent worden bijgehouden.

Vastgestelde knelpunten:
\begin{itemize}
  \item handmatige CSV-imports zijn tijdrovend en foutgevoelig;
  \item vervangingen van toestellen worden niet altijd bijgewerkt;
  \item kabel- en topologierelaties zijn vaak onvolledig of foutief.
\end{itemize}

%------------------------------------------------------------------------------
\section{Onderzoeksvraag}

Hoe kan automatische netwerkdetectie bijdragen aan het up-to-date houden van netwerkdocumentatie en aan een betrouwbare \emph{single source of truth} op geselecteerde WiFi-/AP-sites bij Citymesh, met focus op NetBox-inventarisdata en LLDP-gebaseerde kabel- en topologierelaties?

%------------------------------------------------------------------------------
\section{Aanpak: hybride pipeline}

De oplossing combineert twee complementaire lagen:

\textbf{1. Inventarislaag --- NetBox Discovery met Diode}
\begin{itemize}
  \item Orb agent voert NMAP- en SNMP-/device-discovery uit;
  \item resultaten gaan via gRPC naar Diode, dat ze valideert en in NetBox laadt;
  \item levert betrouwbaar devices, types en interfaces op.
\end{itemize}

\textbf{2. Topologie- en vendorlaag --- \texttt{NetworkTool} (eigen ontwikkeling)}\\
De standaard keten zet LLDP-buren niet automatisch om in kabelobjecten. \texttt{NetworkTool} vult dit aan:
\begin{itemize}
  \item haalt via NAPALM de live LLDP-tabel op en maakt/corrigeert kabelrelaties via de NetBox-API;
  \item drift-detectie herstelt fout gekoppelde kabels;
  \item dry-run vóór elke live synchronisatie;
  \item RUCKUS vSZ-API koppeling voor draadloze access points.
\end{itemize}

\begin{center}
  \captionsetup{type=figure}
  \includegraphics[width=1.0\linewidth]{networktool-network-scanner.png}
  \captionof{figure}{Network Scanner-tab van \texttt{NetworkTool}: keuze tussen een NetBox-gestuurde selectie of handmatige bulkinvoer, met vendor-selectie die de juiste NAPALM-driver activeert.}
\end{center}

%------------------------------------------------------------------------------
\section{Resultaten}

De casestudy omvat vier sites. Per site werd een baseline opgesteld en vergeleken met de toestand na discovery + \texttt{NetworkTool}.

\begin{center}
  \captionsetup{type=table}
  \small
  \begin{tabular}{@{}p{0.27\linewidth} c p{0.16\linewidth} p{0.16\linewidth} p{0.20\linewidth}@{}}
    \toprule
    \textbf{Site} & \textbf{Schaal} & \textbf{Naam\-geving} & \textbf{Type / inventaris} & \textbf{Kabel\-relaties} \\
    \midrule
    Ridefort & 8 & 3/6 $\rightarrow$ 6/6 & 5/6 $\rightarrow$ 6/6 & 3/6 $\rightarrow$ 6/6 \\
    Deleers Go4tech & 61 & 61/61 & 8/11 $\rightarrow$ 11/11 & geen drift \\
    Koninklijke Stallingen & 8 & 8/8 & 5/5 (AP) & 0/7 $\rightarrow$ 7/7 \\
    WZC Riethove & 75 & 75/75 & n.v.t. & 0/74 $\rightarrow$ 57/74\,$^{*}$ \\
    \bottomrule
  \end{tabular}
  \captionof{table}{Voor/na-synthese over de vier sites (voor $\rightarrow$ na). $^{*}$\,17 AP-relaties hingen vast aan één fysiek offline switch en worden bij een volgende discovery aangevuld.}
\end{center}

Daarnaast haalt de vSZ-koppeling per RUCKUS-AP zes velden automatisch op (naam, serienummer, MAC-adres, fabrikant, type en IP-adres).

\textbf{Kerninzicht:} discovery corrigeert niet alleen historische fouten, maar werkt ook als doorlopende kwaliteitscontrole. Op Deleers Go4tech bevatte zelfs documentatie van slechts twee maanden oud nog drie switches met een foutief device type.

%------------------------------------------------------------------------------
\section{Conclusie}

NetBox Discovery met Diode is een sterke en haalbare basis om inventarisdata in NetBox actueel te houden, maar volstaat op zich niet voor een volledig topologische SSOT op WiFi-/AP-sites. De aanvullende \texttt{NetworkTool}-laag voor LLDP-kabelrelaties en de RUCKUS vSZ-koppeling is daarvoor noodzakelijk. Voor alle bereikbare apparatuur werd de documentatie volledig hersteld; de enige resterende lacune is herleidbaar tot één fysiek offline switch.

%------------------------------------------------------------------------------
\section{Aanbevelingen}

\begin{itemize}
  \item zet periodieke discovery in als standaard kwaliteitscontrole;
  \item gebruik steeds een dry-run vóór elke live synchronisatie;
  \item veranker cleanup en naming governance in operationele processen;
  \item neem de RUCKUS vSZ-synchronisatie structureel op;
  \item gebruik een centrale secret manager bij uitrol op tientallen sites.
\end{itemize}

%------------------------------------------------------------------------------
\section{Toekomstig werk}

\begin{itemize}
  \item evaluatie op meer sites en bredere vendormix (bv.\ Cisco, Aruba);
  \item opvolgen of NetBox Labs LLDP-kabels native afdekt in een volgende Diode-/Orb-versie;
  \item integratie met formeel change management voor grootschalige uitrol.
\end{itemize}

\end{multicols}
\end{document}
